\section*{Methods}
\label{sec:methods}

\subsection{Bunch and grid configurations}
\label{sec:configs}

Table~\ref{tab:configs} lists every measured case with the bunch and the
grid, and the section that uses it. The entries are the recorded
configuration values of this work. The generation line for the mid-bunch
bunches is
\seqsplit{generate\_bunch.py --model gamma --gamma G --seed 42 --grid 64 64 32
--cell-size 1e-7 1e-7 1e-7 --centroid 32 32 16 --sigma 6 6 6 --weights
uniform}. The total charge is $N$ times $-e$ for the electron bunches with
one unit weight, where $e = 1.602176634\times10^{-19}$\,C; the kinetic
energy is $(\gamma-1)\,m_ec^2$ with $m_ec^2 = 510.999$\,keV. The mid-bunch
bunches are mono-energetic: the read-back records
$\gamma_{\min}=\gamma_{\max}=1$ for the $\gamma=1$ bunch and
$\gamma_{\min}=\gamma_{\max}=10$ for the $\gamma=10$ bunch.

\begin{table}[t]
\centering
\footnotesize
\setlength{\tabcolsep}{3pt}
\renewcommand{\arraystretch}{1.05}
\caption{Configurations of every measured case and the section that uses it
(recorded values of this work; ``not recorded'' means no value was
recorded). The two mid-slice entries give the bunch generation grid; the
field grid for those cases ($32\times32\times16$ @ $2\times10^{-7}$ m) is
in the second block. Kinetic energy is $(\gamma-1) m_ec^2$; the
multi-velocity driver cell size is $10^{-6}$ m.}
\label{tab:configs}
\begin{longtable}{@{}p{1.9cm}cp{2.2cm}p{2.7cm}p{2.7cm}p{2.7cm}@{}}
\toprule
case & $N$ & total charge ($-e$ each) & size (cells/axis) & spread & used in \\
\midrule
\endfirsthead
\toprule
case & $N$ & total charge ($-e$ each) & size (cells/axis) & spread & used in \\
\midrule
\endhead
\bottomrule
\endfoot
mid slices, g10 & 1000 & $1.602\times10^{-16}$\,C
  & bunch generation grid $64\times64\times32$ @ $10^{-7}$\,m ($6.4\times6.4\times3.2$), $\sigma\!=\!6$ cells/axis
  & none ($\gamma_{\min}\!=\!\gamma_{\max}\!=\!10$)
  & \S\ref{app:slices}, \S\ref{sec:validation} \\
mid slices, g1 & 1000 & $1.602\times10^{-16}$\,C
  & bunch generation grid $64\times64\times32$ @ $10^{-7}$\,m ($6.4\times6.4\times3.2$), $\sigma\!=\!6$ cells/axis
  & none ($\gamma_{\min}\!=\!\gamma_{\max}\!=\!1$)
  & \S\ref{app:slices}, \S\ref{sec:validation} \\
large-domain, g10 & 100 & $1.602\times10^{-17}$\,C
  & extent $2.74\times3.78\times2.65$
  & not recorded
  & \S\ref{app:slices} (Fig.~\ref{fig:large}) \\
large-domain, g1 & 100 & $1.602\times10^{-17}$\,C
  & extent $2.74\times3.78\times2.65$
  & not recorded
  & \S\ref{app:slices} (Fig.~\ref{fig:largeg1}) \\
cost, stationary & $2\times10^6$ & not recorded
  & 20-cell bunch
  & none
  & \S\ref{sec:cost} \\
cost, moving axial & $2\times10^6$ & not recorded
  & 20-cell bunch
  & none
  & \S\ref{sec:cost} \\
landscape, aspect 1 & 1 & $1.602\times10^{-19}$\,C
  & single particle, 96-cell box
  & none
  & \S\ref{sec:gamma-landscape} \\
landscape, aspect 4 & 1 & $1.602\times10^{-19}$\,C
  & single particle at the origin
  & none
  & Table~\ref{tab:d48} \\
 multi-velocity, E 5\% & 200 & $3.2\times10^{-17}$\,C
    & $\sigma\!=\!6$ cells, centroid $(32,32,32)$
    & 5\% energy, $+y$ drift
    & \S\ref{app:multivel} \\
 multi-velocity, A $5^\circ$ & 200 & $3.2\times10^{-17}$\,C
    & $\sigma\!=\!6$ cells, centroid $(32,32,32)$
    & $5^\circ$ angular, $+y$ drift
    & \S\ref{app:multivel} \\
 1\,mrad, binned & 200 & $3.2\times10^{-17}$\,C
    & $\sigma\!=\!6$ cells, centroid $(32,32,32)$, 1.046\,mrad rms $x$-div
    & 1\,mrad, $+y$ drift
    & \S\ref{app:multivel} \\
 divergence, 1\,mrad & 1000 & $1.602\times10^{-16}$\,C
    & $\sigma\!=\!0.6$ each axis, centroid $(3.197,3.208,1.615)$
    & 1\,mrad paraxial, $+y$ (seed 43)
    & \S\ref{app:slices} (reference pending) \\
\end{longtable}

\begin{longtable}{@{}p{2.3cm}p{3.4cm}p{4.5cm}p{4.4cm}@{}}
\toprule
case & $\gamma$, kinetic E ($m_ec^2\!=\!510.999$\,keV) & weight, cell & field grid (domain), host \\
\midrule
\endfirsthead
\toprule
case & $\gamma$, kinetic E ($m_ec^2\!=\!510.999$\,keV) & weight, cell & field grid (domain), host \\
\midrule
\endhead
\bottomrule
\endfoot
mid slices, both & 10 / 5.06\,MeV; 1 / 0 & 1 each; $2\times10^{-7}$\,m & $32\times32\times16$ ($6.4\times6.4\times3.2$), rosi\_v100 \\
large-domain, both & 10 / 5.06\,MeV; 1 / 0 & not recorded; $2\times10^{-7}$\,m & $128\times128\times64$ ($25.6\times25.6\times12.8$), jureca-h100 \\
cost, both & 1 / 0; 10 / 5.06\,MeV & not recorded; $2.5\times10^{-7}$\,m & $66\times68\times66$, jureca-h100 \\
landscape, aspect 1 & 10/50/100 & 1; $10^{-7}$\,m & $96^3$ ($h\!=\!1$), $384^3$ ($h\!=\!1/4$), jureca-h100 (ref: rosi\_v100) \\
landscape, aspect 4 & 1/10/100 & 1; anisotropic, aspect 4 & axis-aligned, $h\!=\!1$, host not recorded \\
 multi-velocity, both & 10 / 5.06\,MeV & 1 each; $10^{-7}$\,m & $64^3$, hal (A30, \texttt{sm\_80}) \\
 divergence & 10 / 5.06\,MeV & 1 each; $2\times10^{-7}$\,m & $32\times32\times16$ ($6.4\times6.4\times3.2$), rosi\_v100 \\
\end{longtable}
\end{table}

\subsection{Command and software versions}
\label{sec:commands}

The binary is invoked as
\seqsplit{picBunch --input bunch.h5 --destination <checkpoint> --config params.toml},
with every unit and conversion parameter in the TOML file. The
\texttt{[picongpu]} section carries \texttt{unit\_length},
\texttt{unit\_e}, \texttt{unit\_b}, \texttt{dt}, \texttt{step},
\texttt{precision}, \texttt{species}; \texttt{[domain]} carries
\texttt{cells}, \texttt{cell\_size}, \texttt{origin};
\texttt{[placement]} carries \texttt{offset\_unit} and the
\texttt{left}/\texttt{front}/\texttt{top} fractions; \texttt{[accuracy]}
carries \texttt{tolerance} and the near-shell settings
(\texttt{near\_shell} \texttt{none}/\texttt{exact}/\texttt{fast}/
\texttt{fast-gpu}, plus \texttt{near\_shell\_radius} and
\texttt{near\_shell\_tile}); \texttt{patch\_count} in \texttt{[placement]}
splits the species into one particle patch per output rank for multi-rank
restarts. A missing conversion key causes a clear error at start-up. The
input bunch is centred by its charge-weighted centroid and placed at the
requested offset from the domain origin; a placement leaving the bunch
outside the domain is rejected at start-up.

The measurement hosts and toolchains are: \texttt{rosi\_v100} (node
\texttt{gv002}, V100 \texttt{sm\_70}, CUDA 12.6,
\texttt{PICBUNCH\_ARCH=70}, Release); \texttt{jrc0880} (H100
\texttt{sm\_90}, CUDA 13.4); \texttt{hal} (NVIDIA A30, \texttt{sm\_80},
CUDA 13.4.92). The openPMD standard is version 5.0~\citep{openpmd}; the
reference PIConGPU is the release whose \texttt{referenceOnly}
field-initializer mode is used in
Section~\ref{sec:validation}. The FFT library is cuFFT in the unnormalized
form of the vendor documentation~\citep{cudacc2024}.

\paragraph{Chunked particle input.}
\texttt{--max-particle-chunk-gb} (a float, for example $1.5$) streams the
selected species in chunks of at most that many gigabytes instead of loading
it once.  Each chunk is read on rank 0, broadcast, computed on the same fixed
grid, and added to the output field.  The mode requires explicit
\texttt{--grid}, \texttt{--origin} and \texttt{--cell-size} and a single
velocity class; the global-domain writer is not combined with it yet.  On a
200-electron test, 1-particle chunks reproduce the unchunked field to
$4.8\times10^{-16}$.

\subsection{Data and code availability}
\label{sec:availability}

The source is available at
\url{https://github.com/psychocoderHPC/picBunch}. The measurement data
referenced by name in this paper are the files recorded in
\texttt{docs/measurements/} of the repository: \texttt{matrix.csv} (the
bench sweep behind Table~\ref{tab:matrix}), \texttt{w4b\_rel\_error.csv}
(Fig.~\ref{fig:Fw4b}, Table~\ref{tab:w4b}), \texttt{coulomb.csv} (the
Coulomb anchor), and the
per-case \texttt{accuracy\_results.json} and \texttt{n1b\_padding\_*.csv}
files (Appendix~\ref{app:padding}), with SHA-256 checksums in
\texttt{MANIFEST.sha256}. The field-slice \texttt{.npz} arrays behind
Figs.~\ref{fig:valforwardE}--\ref{fig:val1mrad} and
Figs.~\ref{fig:large}--\ref{fig:largeg1} are recorded in the measurement
campaign (Appendix~\ref{app:slices}); individual measurements name the
artifact they used. The figures are rendered at their print width from the
recorded data with the shared style module described in the repository
documentation; every figure PDF in \texttt{figs/} was regenerated from
recorded inputs, and the figure audit records the minimum printed font size
and colour-bar label of each.
